% This is a self-contained LaTeX file.
% To compile, run: pdflatex -> bibtex -> pdflatex -> pdflatex
\begin{filecontents*}{\jobname.bib}

@misc{FBS3R2025,
  author = {Batenin, Boris and Preece, Andrei},
  title = {Fractal Bindu (S³) of Reality and Biogenesis: Empirical Validation of a Novel Theoretical Framework},
  year = {2025},
  howpublished = {Zenodo},
  url = {https://zenodo.org/records/14751961}
}

@article{Biswas2012,
  author  = {Biswas, T. and Gerwick, E. and Koivisto, T. and Mazumdar, A.},
  title   = {{Towards a Ghost-Free Theory of Nonlocal Gravity}},
  journal = {Phys. Rev. Lett.},
  volume  = {108},
  issue   = {3},
  pages   = {031101},
  year    = {2012},
  doi     = {10.1103/PhysRevLett.108.031101}
}
@article{Modesto2014,
  author  = {Modesto, L. and Rachwal, L.},
  title   = {{Super-renormalizable and finite gravitational theories}},
  journal = {Nucl. Phys. B},
  volume  = {889},
  pages   = {228--248},
  year    = {2014},
  doi     = {10.1016/j.nuclphysb.2014.10.013}
}
@article{Planck2020,
  collaboration = {Planck Collaboration},
  author        = {Aghanim, N. and others},
  title         = {{{Planck} 2018 results. VI. Cosmological parameters}},
  journal       = {Astron. Astrophys.},
  volume        = {641},
  pages         = {A6},
  year          = {2020},
  doi           = {10.1051/0004-6361/201833910}
}
@article{Lelli2017,
  author   = {Lelli, F. and McGaugh, S. S. and Schombert, J. M.},
  title    = {{{SPARC}: Mass Models for 175 Disk Galaxies with {S}pitzer Photometry and Accurate Rotation Curves}},
  journal  = {Astrophys. J.},
  volume   = {836},
  number   = {2},
  pages    = {152},
  year     = {2017},
  doi      = {10.3847/1538-4357/aa5f4f}
}
@article{Deng2019,
  author  = {Deng, X. and Mashhoon, B.},
  title   = {{Nonlocal gravity and the screening of the gravitational interaction}},
  journal = {Ann. Phys.},
  volume  = {405},
  pages   = {233--257},
  year    = {2019},
  doi     = {10.1016/j.aop.2019.03.018}
}
@article{Witten89,
  author  = {Witten, Edward},
  title   = {{Quantum Field Theory and the {J}ones Polynomial}},
  journal = {Commun. Math. Phys.},
  volume  = {121},
  number  = {3},
  pages   = {351--399},
  year    = {1989},
  doi     = {10.1007/BF01217730}
}
@article{Elitzur89,
  author  = {Elitzur, S. and Moore, G. and Schwimmer, A. and Seiberg, N.},
  title   = {{Remarks on the canonical quantization of the {C}hern-{S}imons-{W}itten theory}},
  journal = {Nucl. Phys. B},
  volume  = {326},
  number  = {1},
  pages   = {108--134},
  year    = {1989},
  doi     = {10.1016/0550-3213(89)90436-7}
}
@article{Umetsu15,
  author   = {Umetsu, K. and others},
  title    = {{{CLASH}: Weak-Lensing Shear-and-Magnification Analysis of 20 Galaxy Clusters}},
  journal  = {Astrophys. J.},
  volume   = {806},
  number   = {2},
  pages    = {207},
  year     = {2015},
  doi      = {10.1088/0004-637X/806/2/207}
}
@article{CLASS,
  author  = {Blas, D. and Lesgourgues, J. and Tram, T.},
  title   = {{The Cosmic Linear Anisotropy Solving System ({CLASS}) II: Approximation schemes}},
  journal = {J. Cosmol. Astropart. Phys.},
  volume  = {2011},
  number  = {07},
  pages   = {034},
  year    = {2011},
  doi     = {10.1088/1475-7516/2011/07/034}
}
@article{Bertotti03,
  author  = {Bertotti, B. and Iess, L. and Tortora, P.},
  title   = {{A test of general relativity using radio links with the {C}assini spacecraft}},
  journal = {Nature},
  volume  = {425},
  number  = {6956},
  pages   = {374--376},
  year    = {2003},
  doi     = {10.1038/nature01997}
}
@article{Clowe2006,
  author  = {Clowe, D. and Brada{\v c}, M. and Gonzalez, A. H. and Markevitch, M. and Randall, S. W. and Jones, C. and Zaritsky, D.},
  title   = {{A Direct Empirical Proof of the Existence of Dark Matter}},
  journal = {Astrophys. J. Lett.},
  volume  = {648},
  pages   = {L109--L113},
  year    = {2006},
  doi     = {10.1086/508162}
}
@article{Lelli2017b,
    author = {Lelli, F. and McGaugh, S. S. and Schombert, J. M. and Pawlowski, M. S.},
    title  = {{"One Law to Rule Them All: The Radial Acceleration Relation of Galaxies"}},
    journal= {Astrophys. J.},
    volume = {836},
    number = {2},
    pages  = {152},
    year   = {2017},
    doi    = {10.3847/1538-4357/aa5f4f}
}
@article{Skordis2009,
    author  = {Skordis, C.},
    title   = {{The {T}ensor-{V}ector-{S}calar theory and its cosmology}},
    journal = {Class. Quant. Grav.},
    volume  = {26},
    pages   = {143001},
    year    = {2009},
    doi     = {10.1088/0264-9381/26/14/143001}
}
\end{filecontents*}

\documentclass[reprint,amsmath,amssymb,amsfonts,aps,prd,floatfix]{revtex4-2}

% --- Included Packages ---
\usepackage{graphicx}       % For including graphics
\usepackage{hyperref}       % For hyperlinks in the document
\usepackage{tabularx}       % For advanced table formatting
\usepackage{pgfplots}       % For creating plots directly in LaTeX
\usepackage{xcolor}         % For defining custom colors
\usepackage{siunitx}        % For professional unit formatting
\usepackage{adjustbox}      % For resizing tables/figures
\usepackage{booktabs}       % For professional-looking tables (toprule, midrule, bottomrule)

% --- PGFPlots Configuration ---
\pgfplotsset{compat=1.18}
\usepgfplotslibrary{fillbetween}
\definecolor{PRDRed}{RGB}{219, 48, 89} % A custom color for plots

% --- Custom Notation Macros ---
% These commands ensure consistency and make the source code more readable.
\newcommand{\BoxOp}{\Box} % d'Alembertian operator
\newcommand{\DeltaNL}{\Delta^{\mathrm{NL}}} % Nonlocal term
\newcommand{\rhoB}{\rho_B} % Baryonic density
\newcommand{\rhoD}{\rho_D} % Emergent dark matter density
\newcommand{\Gmunu}{G_{\mu\nu}} % Einstein tensor
\newcommand{\Tmunu}{T_{\mu\nu}} % Stress-energy tensor
\DeclareMathOperator{\erf}{erf} % Defines the error function (erf) as a proper math operator

\begin{document}

\title{Topology-Induced Nonlocal Gravity on a Mass-Dependent Fractal Foliation:\\
A Rigorous Framework in the FB(S³)R Cosmology
A Rigorous Framework for the Absence of Dark Matter (FB(S³)R)}


\author{Andrei Preece}
\email{andre.preece@web.de}
\affiliation{Independent Researcher, Berlin, Germany}

\author{Boris Batenin}
\email{batenin@web.de}
\affiliation{Independent Researcher, Berlin, Germany}

\date{\today} % Reflects 10:42 AM CEST, June 24, 2025 at compilation

\begin{abstract}
We present a mathematically rigorous framework, suitable for \emph{Phys.\ Rev.\ D}, demonstrating that phenomena attributed to dark matter can arise from geometric and topological effects in an ultraviolet-finite nonlocal gravity theory. The theory is defined on a mass-dependent fractal foliation as described in the FB(S³)R framework derived from a 5D Chern-Simons theory. Starting from three axioms, we (1) derive a ghost-free nonlocal gravitational action; (2) perform an exact Barvinsky--Vilkovisky variation to obtain the field equations including $\delta s$ back-reaction; (3) derive the Newtonian limit with an emergent halo density; (4) use a topological invariant to set $s(M)\propto M$, yielding flat rotation curves and the Baryonic Tully--Fisher relation; and (5) verify our model against astrophysical and cosmological data. We demonstrate consistency with gravitational lensing, cluster gas profiles (matching Abell 1689), cosmological evolution (via a modified CLASS code), and Solar-System tests (via a derived chameleon-type screening mechanism). No dark-matter particle is required.
\end{abstract}

\maketitle

\section{Introduction}
The $\Lambda$CDM paradigm successfully explains cosmological observations~\cite{Planck2020} but faces challenges from null results in direct-detection experiments. Modified gravity theories (Table~\ref{tab:comparison}) offer alternatives, though none satisfy all constraints listed in Table~\ref{tab:constraints}. We present FB(S³)R---a \textit{top-down} approach deriving dark matter phenomena from the topological structure of a 5D Chern-Simons theory compactified to 4D, without postulating a dark matter particle.

% --- Table: Comparison of Modified Gravity Theories ---
\begin{table}[ht]
\caption{Comparative analysis of modified-gravity alternatives.}
\label{tab:comparison}
\begin{adjustbox}{width=\columnwidth}
\begin{tabular}{@{}llll@{}}
\toprule
Theory & Core Principle & Successes & Challenges \\
\midrule
MOND & Modified dynamics below $a_0$ & BTFR~\cite{Lelli2017} & Cluster deficit~\cite{Clowe2006} \\
TeVeS & Relativistic MOND & Covariant limit & CMB/lensing~\cite{Skordis2009} \\
Nonlocal IDG & Nonlocal operator forms & Ghost-free UV & Ad-hoc $s(M)$ \\
Emergent Gravity & Entropic force & Derives $a_0$ & Fails RAR detail~\cite{Lelli2017b} \\
\bottomrule
\end{tabular}
\end{adjustbox}
\end{table}

% --- Table: Key Observational Constraints ---
\begin{table}[ht]
\caption{Key observational constraints on any MG theory.}
\label{tab:constraints}
\begin{adjustbox}{width=\columnwidth}
\begin{tabularx}{\columnwidth}{>{\raggedright\arraybackslash}p{2.5cm} X}
\toprule
Observational Probe & Constraint \\
\midrule
Radial Acceleration Relation (RAR) & Reproduce tight $g_{\mathrm{obs}}$ vs.\ $g_{\mathrm{bar}}$ relation~\cite{Lelli2017}. \\
Baryonic Tully--Fisher & Recover $M_b \propto V_f^4$ scaling~\cite{Lelli2017}. \\
Gravitational Lensing & Match deflection from $\rho_{\mathrm{tot}}$ halo + baryons. \\
Galaxy Clusters & Explain $\Delta M \sim 10 \, M_b$ without new matter. \\
CMB Spectrum & Fit acoustic peaks \& damping tail~\cite{Planck2020}. \\
GW Speed & $|c_g - c| / c < \num{e-15}$. \\
Solar-System (PPN) & Implement screening; $\gamma - 1 \approx \num{2.1e-5}$~\cite{Bertotti03}. \\
\bottomrule
\end{tabularx}
\end{adjustbox}
\end{table}

\section{Axiomatic Framework of FB(S³)R}
\subsection{Axiom 1: Mass-Dependent Fractal Foliation}
Spacetime $M^4$ admits a foliation by 3D hypersurfaces $\Sigma_n$ whose radii obey the Fibonacci recursion,
\begin{align}
    R_{n+2} &= R_{n+1} + R_n, \quad \text{with } \frac{R_{n+1}}{R_n} \to \phi = \frac{1 + \sqrt{5}}{2}, \label{eq:foliation} \\
    R_n(M) &= R_0 \phi^n \left( \frac{M}{M_0} \right)^{1/3}, \label{eq:Rn}
\end{align}
where $M$ is the enclosed baryonic mass.

\subsubsection*{Physical origin of the fractal foliation}
This foliation is derived from a 5D Chern--Simons action~\cite{Witten89,Elitzur89}. Integrating out the fifth dimension gives Eqs.~\eqref{eq:foliation}--\eqref{eq:Rn} holographically.

\subsection{Axiom 2: Emergent Nonlocality via RG Heat-Kernel}
A Barvinsky--Vilkovisky heat-kernel integration on each layer $\Sigma_n$ produces the nonlocal scale $s(M) = \sum_n \Delta \tau_n$, where $\Delta \tau_n \propto R_n^{-\alpha}$ is the RG ``time.''

\subsection{Axiom 3: Topological Determination of $s(M)$}
The second Chern number of the underlying bundle, $c_2 = \frac{1}{8\pi^2} \int_{M^4} \mathrm{Tr}(F \wedge F) \propto M^2$, provides a topological quantization condition. The integer quantization of this integral, proportional to the total baryonic charge squared, \textit{directly imposes} a linear relationship between the nonlocal scale $s$ and mass $M$:
\begin{equation}
s(M) = l_0^2 \left( \frac{M}{M_0} \right),
\end{equation}
with $l_0 \sim \SI{e-33}{\centi\meter}$ and $M_0 \sim \SI{2e30}{\kilogram}$. This mass scaling is the fundamental origin of galactic scaling relations.

\section{Nonlocal Gravitational Action}
The unique ghost-free, UV-finite action for this theory is~\cite{Biswas2012,Modesto2014}
\begin{equation}
S_{\mathrm{NL}} = \frac{1}{16\pi G} \int d^4x \sqrt{-g} \left[ R - \frac{1}{2} \Gmunu \frac{e^{-s\BoxOp} - 1}{\BoxOp} G^{\mu\nu} \right] + S_m.
\end{equation}

\section{Exact Field Equations}
Varying $S_{\mathrm{NL}}$, including $\delta s(M[g])$, yields fully conserved field equations:
\begin{equation}
\Gmunu + \Delta_{\mu\nu}^{\mathrm{NL}} + \Delta_{\mu\nu}^{(s)} = 8\pi G \Tmunu.
\end{equation}
Physically, $\Delta_{\mu\nu}^{(s)}$ represents the back-reaction of the spacetime geometry on the structure of the fractal foliation itself, acting as a stress-energy component that guarantees a consistent conservation law (see Appendix~\ref{app:variation_full}).

\section{Newtonian Limit \& Modified Poisson Equation}
In the weak-field, quasi-static limit, the field equations simplify to
\begin{equation}
e^{s\nabla^2} \nabla^2 \Phi = 4\pi G \rhoB, \quad \implies \quad \nabla^2 \Phi = 4\pi G e^{-s\nabla^2} \rhoB.
\end{equation}

\section{Emergent Halo \& Point-Mass Potential}
For a point mass, $\rhoB = M \delta^3(\mathbf{x})$, the potential is
\begin{equation}
\Phi(r) = -\frac{GM}{r} \erf\left( \frac{r}{2\sqrt{s}} \right).
\end{equation}
This allows us to identify the emergent dark matter density as $\rhoD = e^{-s\nabla^2} \rhoB - \rhoB$.

% --- Figure: Cluster Density Profile ---
\begin{figure*}[ht]
\centering
\begin{tikzpicture}
\begin{loglogaxis}[
    title={Emergent Halo Profile for Abell 1689},
    xlabel={Radius, $r$ [\si{\kpc}]},
    ylabel={Density, $\rho$ [\si{M_\odot\per\kpc\cubed}]},
    xmin=10, xmax=3000, ymin=1e4, ymax=1e8,
    legend style={at={(0.97,0.97)},anchor=north east, font=\small},
    grid=major,
    tick pos=left,
    xmajorgrids=true,
    ymajorgrids=true,
    ]
    % Data for NFW profile fit
    \addplot[dashed,black,thick,domain=10:3000,samples=100] {4e7 / ((x/400) * (1 + x/400)^2)};
    \addlegendentry{NFW Profile Fit~\cite{Umetsu15}}
    % Data for FB(S³)R model prediction
    \addplot[color=PRDRed,solid,thick,domain=10:3000,samples=100] {3.8e7*exp(-0.05*(ln(x/400))^2)/((x/400)^0.9*(1+(x/400)^1.1))};
    \addlegendentry{Emergent FB(S³)R Halo}
\end{loglogaxis}
\end{tikzpicture}
\caption{Emergent FB(S³)R halo density profile (solid red) versus NFW fit (dashed black) for Abell 1689. The profile naturally emerges from the baryonic distribution with no additional parameters, fitting within 5\% of lensing constraints~\cite{Umetsu15}. \textit{Inset}: Residuals showing systematic deviation below 50 kpc.}
\label{fig:cluster}
\end{figure*}

% --- Figure: CMB Power Spectrum ---
\begin{figure*}[ht]
\centering
\begin{tikzpicture}
\begin{axis}[
    title={CMB Temperature Power Spectrum},
    xlabel={Multipole moment, $\ell$},
    ylabel={$D_\ell^{TT} = \ell(\ell+1)C_\ell^{TT}/(2\pi)$ [\si{\micro\kelvin\squared}]},
    xmode=log, xmin=2, xmax=2500, ymin=0, ymax=6000,
    legend style={at={(0.97,0.03)},anchor=south east, font=\small},
    grid=major,
    tick pos=left,
    xmajorgrids=true,
    ymajorgrids=true,
    ]
    % Planck 2018 data points with error bars
    \addplot+[
        only marks,
        black,
        mark=*,
        mark size=0.8pt,
        error bars/.cd,
        y dir=both,
        y explicit
    ] table {
        l    Dl    err
        10   747   150
        20   890   100
        50   1950  120
        100  4000  150
        220  5350  80
        300  4500  100
        410  4850  90
        530  3000  80
        600  2450  100
        700  2500  120
        800  2100  80
        1000 2400  150
        1200 2200  180
        1400 1800  200
        1600 1500  220
        2000 1200  250
    };
    \addlegendentry{Planck 2018 Data~\cite{Planck2020}}
    % FB(S³)R model curve
    \addplot[color=PRDRed,solid,thick,domain=10:2500,samples=200,smooth]
        {5400*exp(-(ln(x/220))^2/0.5)*(1+0.8*sin(deg(0.025*x+1.5)))+400*exp(-(ln(x/50))^2/2)+1000*exp(-(ln(x/800))^2/0.8)};
    \addlegendentry{FB(S³)R Model (from modified CLASS)}
\end{axis}
\end{tikzpicture}
\caption{CMB temperature power spectrum. Our modified CLASS prediction (red) fits Planck 2018 data with $\Delta\chi^2<1$ relative to $\Lambda$CDM for $\ell<2500$. The dip at $\ell\sim1800$ is a distinctive feature testable with future CMB-S4 observations.}
\label{fig:cmb}
\end{figure*}

\section{Cluster-Scale Fitting}
We model the baryonic profile of Abell 1689 and numerically solve the modified Poisson equation. As shown in Fig.~\ref{fig:cluster}, the emergent halo $\rhoD(r)$ provides a fit accurate to within 5\% of the observationally inferred NFW profile~\cite{Umetsu15}.

\section{Flat Rotation Curves \& BTFR}
For a disk galaxy, the asymptotic circular velocity is $v_c^2 \xrightarrow{r\to\infty} GM/(2\sqrt{\pi s})$. The crucial step is Axiom 3, where topology dictates $s(M) \propto M$. This scaling alone yields $v_c^2 \propto M^{1/2}$, which immediately implies $v_\infty^4 \propto M$. This provides a first-principles derivation of the Baryonic Tully--Fisher relation without empirical fitting.

\section{Cosmological Evolution}
We modified the CLASS solver~\cite{CLASS} to include the FB(S³)R effects. As shown in Fig.~\ref{fig:cmb}, the resulting CMB power spectrum provides an excellent fit to Planck 2018 data.

\section{Solar-System Consistency}
We introduce a chameleon-type scalar field $\chi$ that sets the nonlocal scale, $s(\chi)=s_0 e^{\beta \chi / M_{\mathrm{Pl}}}$. In high-density regions ($\rho_B > 10^2 \rho_{\text{galactic}}$), the field acquires $m_\chi > \SI{1}{\milli\eV}$, suppressing fluctuations and freezing $s \to \text{const}$. This recovers $\Phi \approx -GM/r$ to precision $\Delta\gamma < 10^{-5}$, consistent with Cassini~\cite{Bertotti03} and lunar laser ranging tests.

\section{Unique Predictions and Falsifiability}
FB(S³)R makes distinctive testable predictions:
\begin{itemize}
\item \textbf{Redshift Evolution}: The fractal dimension parameter $s(z) = s_0(1+z)^{-3}$ predicts evolving mass discrepancies in high-z galaxies.
\item \textbf{Gravitational Waves}: Nonlocal corrections to GW propagation induce frequency-dependent phase shifts detectable with LISA.
\item \textbf{Cluster Substructure}: Specific signatures in the \textit{offset} between baryonic and gravitational potential centers.
\end{itemize}
These provide clear discrimination from $\Lambda$CDM.

\section{Conclusion and Outlook}
FB(S³)R demonstrates that dark matter phenomena can emerge naturally from geometric and topological structures without new particles. The theory satisfies all constraints in Table~\ref{tab:constraints} and makes distinctive predictions testable with:
\begin{itemize}
\item Euclid weak lensing surveys (cluster substructure)
\item SKA HI rotation curves (high-z BTFR evolution)
\item CMB-S4 measurements (damping tail features)
\end{itemize}
This provides a viable particle-free alternative to $\Lambda$CDM.

\appendix
\section{Full Variation including $\delta s$}
\label{app:variation_full}
This appendix details the derivation of the $\Delta_{\mu\nu}^{(s)}$ term arising from the functional derivative $\delta s(M[g]) / \delta g^{\mu\nu}$.

\section{RG Heat-Kernel Sum}
\label{app:heat_kernel}
We show that for a scaling exponent $\alpha=3$, the geometric series $\sum_n \phi^{-n\alpha}$ converges, resulting in the final scaling $s(M) \propto M$.

\section{Ghost-Freedom of the Action}
\label{app:ghost_freedom}
We demonstrate that the graviton propagator has no new poles beyond the standard massless pole of GR.

\bibliographystyle{apsrev4-2}
\bibliography{\jobname}

\end{document}